% chapter_09_G4.tex -- chapter 09 of the JitJet book: the hand-over of the generator final state to Geant4.
% Every number comes from test/09_G4/out/ (ledger_09_G4.txt, extras_09_G4.txt, simAnalyzer_signal_09_G4.log,
% simAnalyzer_pileup_09_G4.log, step2_G4_signal_09_G4.log, the PNG pages), from the FWLite extras and the release files
% listed in doc/chapter/chapter_09_G4.md, section 4.

\chapter{09 -- G4: the 634 stable particles are handed to Geant4, 571 are accepted}
\label{ch:09_G4}

\section{What this chapter is about}

Everything up to here was Pythia8: a list of particles with four-momenta and production vertices, but no
detector. From this chapter on the list is given to Geant4 (``GEometry ANd Tracking''), the toolkit with which
the LHC experiments simulate how particles cross matter. Geant4 needs three things. A \emph{geometry}: the
detector as a tree of volumes, each with a shape and a material. A \emph{physics list}: the processes a particle
can undergo in a material (a decay, an ionisation, a photon conversion, a nuclear interaction), with their cross
sections. And \emph{user actions}, through which the experiment says which particles to track, when to stop them
and what to record. Geant4 then moves every particle through the geometry step by step, and the volumes declared
\emph{sensitive} write down what crosses them. In CMSSW this is the \code{SIM} step (\code{cmsDriver.py -s SIM}),
one module, \code{g4SimHits} (\code{OscarMTProducer}, configured by
\path{SimG4Core/Application/python/g4SimHits_cfi.py}).

Its first job is the \emph{hand-over}: it reads the HepMC record \code{generatorSmeared}, decides which of the
stable particles are worth simulating, turns them into Geant4 \emph{primaries} and only then starts tracking
them. This chapter is about that decision alone. To see it without anything else in the way, the script runs
Geant4 with its physics reduced to the transportation: no magnetic field, no decays in flight, no electromagnetic
and no hadronic processes (the physics list \code{JourneyTransportPhysics} of
\path{JitJetSW/Journey/plugins/JourneyTransportPhysics.cc}, which constructs every Geant4 particle and registers no
process). Every accepted particle then flies in a straight line, loses no energy, and the file shows exactly which
particles Geant4 accepted, which it dropped and where each one would go if nothing happened to it: every number of
this chapter is a statement about the cuts and the geometry, and every difference to chapter~\ref{ch:10_SIM} will
be a statement about the physics (the run is also fast, about a minute for the 250 library collisions). The same is done
for the 250 library collisions of chapter~\ref{ch:08_PU}, so that the ledger can be closed for the signal and for
its 68 in-time pileup collisions. The full physics is switched on, piece by piece, in chapter~\ref{ch:10_SIM}.
Figure~\ref{fig:09_chain} is the chain the particles go through, and table~\ref{tab:09_glance} what each link of
it does to our event.

\begin{figure}[!htb]
\centering
\begin{tikzpicture}[x=0.93cm, font=\footnotesize, >={Stealth[length=2mm]},
  box/.style={draw, rounded corners=2pt, minimum height=10mm, text width=3.25cm, align=center, inner sep=2.5pt},
  gen/.style={box, fill=yellow!18}, hand/.style={box, fill=orange!15}, g4/.style={box, fill=blue!10},
  prod/.style={box, fill=green!12}, geo/.style={box, fill=black!6, text width=6.2cm}]
\node[gen] (hep) at (0,1.5) {the generator event\\{\tiny\ttfamily generatorSmeared (HepMC)}};
\node[hand] (cut) at (4.25,1.5) {the hand-over cuts\\{\tiny\ttfamily Generator.cc}};
\node[hand] (pri) at (8.5,1.5) {Geant4 primaries\\{\tiny\ttfamily G4PrimaryParticle}};
\node[hand] (stk) at (12.75,1.5) {neutrinos killed\\{\tiny\ttfamily StackingAction}};
\node[g4] (trk) at (12.75,-0.5) {tracking: transportation only\\{\tiny\ttfamily JourneyTransportPhysics}};
\node[g4] (sd) at (8.5,-0.5) {sensitive detectors\\{\tiny\ttfamily TrackerSD, CaloSD, HF library, CaloTrkProcessing}};
\node[prod] (out) at (4.25,-0.5) {the SIM products\\{\tiny\ttfamily SimTrack, SimVertex, PSimHit, PCaloHit}};
\node[geo] (det) at (10.6,-2.35) {the detector: volumes, materials, regions, sensitive volumes\\{\tiny\ttfamily GeometryFileRcd: 141DD4hepV2\_Extended2024}};
\draw[->] (hep) -- (cut); \draw[->] (cut) -- (pri); \draw[->] (pri) -- (stk); \draw[->] (stk) -- (trk);
\draw[->] (trk) -- (sd); \draw[->] (sd) -- (out);
\draw[->] (det.north-|trk) -- (trk.south); \draw[->] (det.north-|sd) -- (sd.south);
\draw[->, dashed] (out.west) -- ++(-1.35,0) node[left, align=center, font=\scriptsize] {chapter\\10};
\node[font=\scriptsize, text=black!60] at (4.25,2.3) {\ref*{sec:09_cuts}, \ref*{sec:09_decays}};
\node[font=\scriptsize, text=black!60] at (12.75,2.3) {\ref*{sec:09_cuts}};
\node[font=\scriptsize, text=black!60, anchor=west] at (12.85,-1.35) {\ref*{sec:09_boundary}};
\node[font=\scriptsize, text=black!60] at (8.5,0.3) {\ref*{sec:09_hits}};
\node[font=\scriptsize, text=black!60, anchor=east] at (det.west) {\ref*{sec:09_detector}};
\end{tikzpicture}
\caption{The \code{SIM} step of this chapter: from the generator record to the simulation products. Orange: the
hand-over, which decides what Geant4 simulates; blue: the tracking and the recording, here with the physics reduced
to the transportation; grey: the detector, read from the global tag. The grey numbers are the subsections of
section~\ref{sec:09_numbers} that explain each box.}
\label{fig:09_chain}
\end{figure}

\begin{table}[!htb]
\centering\small\setlength{\tabcolsep}{4pt}
\caption{The chapter at a glance: one subsection of section~\ref{sec:09_numbers} per ingredient of figure~\ref{fig:09_chain}, what it does and what
it gives for our event.}
\label{tab:09_glance}
\begin{tabular}{@{}c>{\raggedright\arraybackslash}p{3.3cm}>{\raggedright\arraybackslash}p{6.4cm}>{\raggedright\arraybackslash}p{5.0cm}@{}}
\toprule
section & what & what it does & in our event \\
\midrule
\ref{sec:09_detector} & the detector & the 2024 geometry of the global tag: volumes, materials, sensitive volumes, regions with production cuts & 15 kinds of sensitive detector, 32 hit collections \\
\ref{sec:09_cuts} & the hand-over cuts & keeps a stable particle with $p>40\MeV$ whose straight line leaves a $2.9\cm$ cylinder before $|z|=354.8\cm$; kills the neutrinos & 634 stable particles, 571 primaries, 63 dropped \\
\ref{sec:09_decays} & the decays Pythia made & keeps the generator decays inside the beam pipe; the daughters start where the parent decayed & 29 \code{SimVertex}s, the farthest $25.8\cm$ from the collision \\
\ref{sec:09_boundary} & straight lines & every primary flies straight to the boundary between the tracker and the calorimeters & 557 of the 571 reach it \\
\ref{sec:09_hits} & hits without physics & the sensitive detectors mark each crossing; the HF shower library answers by itself & 6734 tracker hits, 1511 HF photoelectrons \\
\ref{sec:09_pileup} & the pileup & the same for the 68 in-time collisions of our bunch crossing & $12\,993$ stable, $10\,388$ primaries \\
\bottomrule
\end{tabular}
\end{table}

\begin{plainwords}[The hand-over]
Geant4 does not read \code{genParticles}; it reads the HepMC event and builds a \code{G4PrimaryParticle} for
every generator particle that passes the cuts of the \code{Generator} parameter set: a momentum above
$40\MeV$, a direction within $\abseta<5.5$, and, for particles the generator has already decayed, a decay vertex
outside a cylinder of radius $2.9\cm$ around the beam (the beam pipe), in which case Geant4 is given the
undecayed parent instead. Neutrinos pass these cuts and are killed by the \code{StackingAction}
(\code{TrackNeutrino = False}) before their first step. Each accepted particle becomes a \code{SimTrack} whose
\code{genpartIndex()} is the HepMC barcode of the generator particle, so the link between the generator record
and the simulation is kept, particle by particle. The dropped ones leave no trace in the simulation at all.
\end{plainwords}

Section~\ref{sec:09_numbers} follows figure~\ref{fig:09_chain} in six subsections, one per ingredient; each
says what the ingredient does, where its parameters are set, and what it gives for our event.

\section{From the previous chapter}

Two inputs enter (Rule E): the signal, \path{test/07_GenJet/out/step2_GENJET_07_GenJet.root}, the full
generator event of chapter~\ref{ch:06_Decay} with its generator jets, 634 stable particles carrying
$13\,600\GeV$ and $528.9\GeV$ of scalar \pt{}; and the pileup library,
\path{test/08_PU/out/step1_GEN_MinBias_08_PU.root}, the 250 minimum-bias collisions of chapter~\ref{ch:08_PU},
of which library events 72 to 139 are the 68 in-time collisions of our bunch crossing (the ones the
\code{MixingModule} of chapter~\ref{ch:10_SIM} draws into bunch crossing 0 for this signal event: seed 12345,
Poisson with mean 60, library events 1--71 going to bunch crossing $-1$), $12\,993$ stable particles with
$924.8\TeV$ and $5315\GeV$ of \pt{}.\footnote{All numbers of this chapter are in
\path{test/09_G4/out/ledger_09_G4.txt} (the output of \code{common/g4Ledger.py}),
\path{extras_09_G4.txt} (the output of \path{test/09_G4/g4Extras_09_G4.py}), \path{simAnalyzer_signal_09_G4.log}
and \path{simAnalyzer_pileup_09_G4.log}; the few extra ones (the dropped particles by kind, the displaced vertices)
were read from the SIM files with FWLite, and the dimensions of the detector from the XML of the release; the
companion lists them all in section~4.} Nothing in the generator record changes: the \code{SIM} process copies
\code{genParticles}, \code{generatorSmeared} and the jets unchanged and adds its own products under the
process name \code{SIM}: \code{SimTrack}s and \code{SimVertex}s, the \code{PSimHit} collections of the tracker
and muon system (one per subdetector and time-of-flight window, e.g.\ \code{TrackerHitsTIBLowTof}) and the
\code{PCaloHit} collections of the calorimeters (\code{EcalHitsEB}, \code{EcalHitsEE}, \code{EcalHitsES},
\code{HcalHits}). The smeared primary vertex of the signal, drawn in chapter~\ref{ch:06_Decay}, is where all
primaries start: $(0.028, -0.017, 3.11)\cm$ for the signal, $(0.029, -0.017, -0.48)\cm$ for library event~72. The
$3.11\cm$ of the signal will matter for the $\eta$ cut (section~\ref{sec:09_cuts}).

The \code{SIM} job runs with the 2025 settings of the release (global tag \code{auto:phase1\_2025\_realistic},
era \code{Run3\_2025}; the signal was generated with those of 2024, which at generator level are the same,
chapter~\ref{ch:08_PU}). The detector given to Geant4 comes from the global tag, and in this release the 2025 tag
has the geometry and the field map of 2024: \code{GeometryFileRcd} \code{Extended} is
\code{141DD4hepV2\_Extended2024} in both, the field \code{MFGeometry\_160812} at 3.8\,T too. None of the five
modifiers that \code{Run3\_2025} adds to \code{Run3\_2024} appears in \code{g4SimHits\_cfi.py}. Rerun with the
2025 settings, both files of this chapter are the 2024 ones, primary for primary and hit for hit (companion,
section~4): the 2025 detector of CMSSW\_15\_0\_5 differs from that of 2024 only in conditions and modifiers
that the simulation does not read.

The two objects to follow are the $b$ and $\bar b$ systems of chapter~\ref{ch:07_GenJet}: the nine stable
decay products of the $B^-$ ($\pt=94.3\GeV$, $\eta=-1.57$, among them the $95.5\GeV$ antineutrino) and the
16 of the $B^{*+}$ ($\pt=114.2\GeV$, $\eta=-3.75$, no neutrino).

\section{The numbers of the event}
\label{sec:09_numbers}

\subsection{The detector that Geant4 is given}
\label{sec:09_detector}

\textbf{What it does.} Geant4 cannot move a particle without a world to move it in. The world of the \code{SIM}
step is read from the global tag: the record \code{GeometryFileRcd}, label \code{Extended}, holds
\code{141DD4hepV2\_Extended2024}, the description of the 2024 detector in the DD4hep format, the same description
that the release carries as XML in \path{Geometry/CMSCommonData/data/dd4hep/cmsExtendedGeometry2024.xml} (a list of
374 included files, a few per subdetector). Each \emph{volume} has a shape, a material and a place inside a
\emph{mother} volume, so the detector is a tree. Its trunk is a chain of air-filled envelopes (\code{World}, \code{OCMS},
\code{CMSE}); inside them five mothers matter for this book (figure~\ref{fig:09_rz}): \code{BEAM}, the beam pipe;
\code{Tracker}, the silicon detectors and their supports, a cylinder of radius $123.3\cm$ and half-length
$293.2\cm$, with a thin ``nose'' around the beam pipe that continues to $|z|=355.1\cm$; \code{CALO}, the
electromagnetic and hadronic calorimeters of the barrel and the endcaps, out to $r=295\cm$; \code{MUON}, the
solenoid, the iron yoke and the muon chambers beyond; and, $11.15$\,m from the collision point on either side,
\code{VCAL}, the forward calorimeter HF, a tube from $r=12.5$ to $159.5\cm$.

Beyond shape and material, the XML attaches labels to volumes (blocks called \code{SpecPar}), and these turn a
geometry into a detector simulation. \emph{Sensitive detectors}: a volume labelled with the sensitive detector
\code{TkAccumulatingSensitiveDetector} and the read-out name \code{TrackerHitsPixelBarrel} records every step
taken inside it into that collection (the files \code{trackersens.xml}, \code{ecalsens.xml},
\code{hf.xml}, \code{muonSens.xml}). \emph{Regions}: groups of volumes with their own production cuts (the box
below). And the \emph{boundary}: the calorimeter mothers carry \code{CaloTrkProcessing}
(\path{Geometry/HcalSimData/data/CaloUtil/2021/v1/CaloUtil.xml}), which notices when a track steps into
\code{CALO}, \code{VCAL} or another calorimeter volume from one of the volumes listed as ``inside'' (\code{BEAM},
\code{Tracker}, \code{MUON}, \ldots) and stores where (section~\ref{sec:09_boundary}).

\begin{figure}[!htb]
\centering
\begin{tikzpicture}[x=0.0104cm, y=0.0135cm, font=\scriptsize, >={Stealth[length=1.6mm]}]
\fill[green!12] (0,123.3) -- (293.2,123.3) -- (293.2,21.5) -- (355.15,21.5) -- (355.15,7.76) -- (554.1,8.93)
  -- (554.1,295) -- (0,295) -- cycle;
\fill[blue!10] (0,0) rectangle (293.2,123.3); \fill[blue!10] (293.2,2.9) rectangle (355.15,21.5);
\fill[red!12] (1114.95,12.495) rectangle (1280.05,159.5);
\fill[orange!10] (0,295) rectangle (554.1,330);
\fill[black!25] (0,0) rectangle (1300,2.9);
\draw[black!50] (0,123.3) -- (293.2,123.3) -- (293.2,21.5) -- (355.15,21.5) -- (355.15,7.76) -- (554.1,8.93)
  -- (554.1,295) -- (0,295);
\draw[black!50] (1114.95,12.495) rectangle (1280.05,159.5);
\node at (140,62) {\code{Tracker}};
\node at (300,210) {\code{CALO}: ECAL + HCAL};
\node at (300,312) {\code{MUON} beyond $r=295\cm$};
\node[align=center] at (1197,95) {HF\\(\code{VCAL})};
\node[anchor=north, text=black!60] at (820,-6) {beam pipe, $r<2.9\cm$};
% the lines from the signal vertex (z = +3.11 cm, drawn at x = -3.11)
\fill (-3.11,0) circle (1.2pt);
\draw[red!75!black, thick] (-3.11,0) -- (662.7,295);
\node[red!75!black, anchor=west] at (705.8,308) {$b$ system, $\eta=-1.55$};
\draw[blue!70!black, thick] (-3.11,0) -- (1280.05,60.57);
\node[blue!70!black, anchor=south] at (790,42) {$\bar b$ system, $\eta=-3.75$};
\draw[black!60, dashed] (-3.11,0) -- (705.8,295);
\node[black!60, anchor=west] at (705.8,280) {corner: $|\eta|=1.61$};
\draw[black!60, dotted, thick] (0,0) -- (1280.05,10.46);
\node[text=black!70, anchor=south] at (920,10) {$|\eta|=5.5$};
\fill[red!75!black] (275.2,123.3) circle (1.6pt); \fill[blue!70!black] (355.15,16.9) circle (1.6pt);
\fill[blue!70!black] (1114.95,52.8) circle (1.6pt);
% axes
\draw[->] (0,0) -- (1330,0) node[right] {$-z$ [cm]};
\draw[->] (0,0) -- (0,345) node[left] {$r$ [cm]};
\foreach \z in {293.2, 355.15, 554.1, 1114.95, 1280.05} \draw (\z,0) -- (\z,-8);
\node[below] at (275,-6) {293}; \node[below] at (375,-6) {355}; \node[below] at (554.1,-6) {554};
\node[below] at (1100,-6) {1115}; \node[below] at (1295,-6) {1280};
\foreach \r in {123.3, 295} \draw (0,\r) -- (-8,\r) node[left] {\r};
\end{tikzpicture}
\caption{The detector of the \code{SIM} step in the $r$--$z$ plane, for $z<0$, where both jets of the event go
(\code{cms.xml} and \code{hcalforwardalgo.xml} of the release;
the $r$ and $z$ scales differ, the muon system beyond $|z|=554\cm$ is not drawn). Blue: the \code{Tracker} volume
with its nose along the beam pipe; green: \code{CALO}; red: the HF. Lines: the straight flights from the signal
vertex at $z=+3.11\cm$ of the $b$ system (red), of the $\bar b$ system (blue), to the corner of the tracker volume
(dashed), and the line $|\eta|=5.5$ from the centre (dotted); the dots are where the systems cross into a
calorimeter volume (section~\ref{sec:09_boundary}).}
\label{fig:09_rz}
\end{figure}

\begin{plainwords}[A production cut is a range, not an energy]
Geant4 does not create every secondary particle a process could make: a photon, electron or positron that would
travel less than the \emph{production cut} of its region before stopping is not created, and its energy is
deposited on the spot. The cut is given as a distance and converted, material by material, into an energy, so
the same $1\mm$ is a larger energy in the dense crystals of the ECAL than in air. Small cuts where the detail
matters (the silicon, which must see the delta electrons that widen a cluster), large ones where it does not
(the iron of the muon system, where only the muons count): the cuts trade precision for time.
\end{plainwords}

\textbf{The parameters.} Which of the labelled sensitive detectors are switched on is decided in python, by
\code{OnlySDs} of \code{g4SimHits\_cfi.py}: for Run~3 (modifier \code{run3\_common}) 15 kinds, writing 25
\code{PSimHit} collections (the six tracker subdetectors, each in a \code{LowTof} and a \code{HighTof} collection,
the five muon technologies, the luminosity and forward detectors) and 7 \code{PCaloHit} collections (ECAL barrel,
endcap and preshower, HCAL, ZDC and two more). The regions and their production cuts are in the XML
(table~\ref{tab:09_prodcuts}); outside any region the cut is \code{DefaultCutValue = 1.0}\,cm
(\code{g4SimHits\_cfi.py}). The dimensions of figure~\ref{fig:09_rz} are constants of
\path{Geometry/CMSCommonData/data/cms/2021/v3/cms.xml}:
\code{TrackCalorR = 1.233\,m}, \code{TrackBeamZ2 = 2.932\,m}, \code{PIXNoseLength = 619.5\,mm},
\code{CalorMuonR = 2.950\,m}, and of \path{Geometry/HcalCommonData/data/hcalforwardalgo.xml}: the HF from $|z|=11.1495$ to $12.8005$\,m.

\begin{table}[!htb]
\centering\small
\caption{The regions of the 2024 geometry and their production cuts for photons, electrons and positrons
(\code{CMSCutsRegion} and \code{ProdCutsFor*} in the \code{*ProdCuts.xml} files included by
\code{cmsExtendedGeometry2024.xml}).}
\label{tab:09_prodcuts}
\setlength{\tabcolsep}{8pt}\begin{tabular}{@{}lll@{}}
\toprule
region & volumes & cut \\
\midrule
\code{BeamPipe} & the beryllium tube at the centre & $0.01\mm$ \\
\code{TrackerSensRegion} & the active silicon of the strip modules & $0.1\mm$ \\
\code{Tracker(Pixel)DeadRegion} & supports, cables, cooling & $1\mm$ \\
\code{EcalRegion} & the crystals and their housing & $1\mm$ \\
\code{APDRegion} & the photodiodes of the barrel ECAL & $0.005\mm$ \\
\code{PreshowerSensRegion} & the silicon of the preshower & $0.1\mm$ \\
\code{HcalRegion} & \code{CALO}, \code{VCAL} & $1\mm$ \\
\code{Muon} & \code{MUON} & $100\mm$ ($\gamma$), $10\mm$ ($e^\pm$) \\
\bottomrule
\end{tabular}
\end{table}

\textbf{The numbers of the event.} In this chapter the regions do nothing: with no process registered no
secondary can be created, whatever its range, and page~8 of the analyzer is empty
(figure~\ref{fig:09_boundary}, right). Of the 32 hit collections the job writes, 7 are filled for the signal: the
six \code{LowTof} tracker collections, by crossings without energy loss, and \code{HcalHits}, by the HF shower
library (section~\ref{sec:09_hits}). The geometry acts through its boundaries alone, and these decide where each
primary is recorded (section~\ref{sec:09_boundary}).

\subsection{The hand-over cuts}
\label{sec:09_cuts}

\textbf{What it does.} The class \code{Generator} (\path{SimG4Core/Generators/src/Generator.cc}) loops over the
vertices of the HepMC event and, for every particle, decides whether Geant4 gets it. A stable particle (status 1)
is tested in a fixed order: if it is a neutron, antineutron or photon so forward that it would stay in the beam
pipe for 52.8\,m, it is kept (the ``very-forward'' rule); otherwise it must have a momentum between $40\MeV$ and
$99\,999\GeV$, and its straight line must leave the beam-pipe cylinder in time (the $\eta$ cut, box below). A
particle that passes becomes a \code{G4PrimaryParticle} attached to a \code{G4PrimaryVertex} at its production
point, with \code{setGenId} set to its HepMC barcode. Particles that the generator already decayed (status 2) are
treated in section~\ref{sec:09_decays}. The \code{StackingAction}, which sees every track before Geant4 starts it,
then kills the neutrinos.

\textbf{The parameters.} Table~\ref{tab:09_cuts} lists the four reasons for which a stable particle gets no
\code{SimTrack}, with the parameter that implements it and its effect on the signal and on the 68 pileup
collisions.

\begin{table}[htb]
\centering\small\setlength{\tabcolsep}{4pt}
\caption{Why a stable particle gets no \code{SimTrack}: the cuts of \code{g4SimHits\_cfi.py} and what they remove
from the signal and from the 68 in-time pileup collisions (\code{ledger\_09\_G4.txt}; the split by kind from the
FWLite extras of the companion). The last line is the exception that keeps very forward neutrals.}
\label{tab:09_cuts}
\begin{tabular}{p{1.9cm}p{3.5cm}p{4.0cm}p{4.7cm}}
\toprule
cut & parameter & signal (634 stable) & 68 collisions ($12\,993$ stable) \\
\midrule
$\abseta>5.5$ & \code{Generator}: \code{ApplyEtaCuts}, \code{MinEtaCut} / \code{MaxEtaCut} $=\mp5.5$; the line inside $r<2.9\cm$ at $|z|=355\cm$ & 56, $6581\GeV$, $19.6\GeV$ \pt{}: 27 charged hadrons (3869, of which 4 protons 1277), 8 neutral hadrons (1580), 21 photons (1132); $5.52<\abseta<8.04$ & 2492, $706.0\TeV$, $851.4\GeV$ \pt{}: 1228 charged hadrons ($517.0\TeV$, 132 protons with $226.2\TeV$), 251 neutral hadrons (98.7), 994 photons (89.4), 19 electrons (0.9); 211 beyond $\abseta=8.2$ ($364.9\TeV$) \\
$p<40\MeV$ & \code{ApplyPCuts}, \code{MinPCut = 0.04} & 6 photons from $\pi^0$, $0.2\GeV$ & 98 (94 photons, 2 charged hadrons, 2 electrons), $2.7\GeV$ \\
neutrino & \code{StackingAction}: \code{TrackNeutrino = False} & 1 ($\bar\nu_e$ of the $B^-$: $95.5\GeV$, $\pt=37.5$) & 7, $57.5\GeV$ \\
other & the $\eta$ line from a displaced vertex & 0 & 8 ($166.4\GeV$): 4 photons, 4 charged hadrons at $\abseta=5.48$--$5.50$ \\
\midrule
kept beyond 5.5 & forward $n$, $\bar n$, $\gamma$ inside $r<2.9\cm$ at $|z|=52.8$\,m & 2: a neutron ($151\GeV$, $\eta=9.72$), a $\pi^+$ ($84.6\GeV$, $\eta=-5.50$) & 168, $134.2\TeV$ \\
\bottomrule
\end{tabular}
\end{table}

\begin{plainwords}[The $\eta$ cut is a line, not an angle]
The cut does not test $\abseta<5.5$. It follows the straight line of the particle from its production vertex to
$|z| = 2.9\cm\times\sinh5.5 = 2.9\times122.3 = 354.8\cm$ and drops the particle if the line is then still inside
the $2.9\cm$ cylinder (\code{RDecLenCut}). From the nominal collision point this is the same as $\abseta>5.5$,
since a line of pseudorapidity $\eta$ is at $r=|z|/\sinh|\eta|$. Our vertex is at $z=+3.11\cm$. The $\pi^+$ of
$84.6\GeV$ at $\eta=-5.5002$ flies towards negative $z$: from the nominal point it would arrive at
$r=354.8/\sinh5.5002=2.899\cm$, $10\um$ inside, and be dropped; from our vertex it flies $357.9\cm$ and arrives at
$2.940\cm$: kept. The photon of $12.9\GeV$ at $\eta=-5.5240$ arrives at $2.887\cm$: dropped. Towards positive $z$
the vertex shortens the flight: the $K^0_S$ at $\eta=+5.4824$ arrives at $2.903\cm$ instead of $2.951\cm$, kept by
$30\um$. The very-forward rule uses the same construction with $\eta=8.2$: $2.9\cm\times\sinh8.2 = 2.9\times1820 =
5279\cm$, the 52.8\,m of the table.
\end{plainwords}

\textbf{The numbers of the event.} The $\eta$ cut is where almost all the energy goes: 56 particles with
$6581\GeV$ of the signal, at $5.52<\abseta<8.04$, and in the pileup eight particles at $\abseta = 5.48$--$5.50$
with vertices displaced by $1$--$6\cm$ in $z$ are dropped although their momentum points just inside $5.5$ (the
``other'' column of the ledger, the box above at work). The very-forward rule
(\code{EtaCutForHector = 8.2}) hands the $151\GeV$ neutron at $\eta = 9.72$ of the signal and 168 particles with
$134.2\TeV$ of the 68 pileup collisions to Geant4 for the Zero Degree Calorimeter; they are later stopped in the
dead regions of the beam line (\code{DeadRegions = QuadRegion, InterimRegion}). The $40\MeV$ cut removes six
photons from $\pi^0$ decays ($27$--$38\MeV$), $0.2\GeV$ in all. The one neutrino is the $\bar\nu_e$ of the
semileptonic $B^-$ decay with $95.5\GeV$ and $37.5\GeV$ of \pt{}: the $b$ system enters Geant4 with 8 of its 9
particles and $56.8$ of its $94.3\GeV$ of \pt{}, while the $\bar b$ system enters complete, 16 of 16,
$114.2\GeV$. The net $b+\bar b$ \pt{}, $20.3\GeV$ at hadron level, is $57.5\GeV$ among the primaries: a detector
cannot see the neutrino, and the imbalance it will measure for this event is decided here.

\subsection{The decays that Pythia already made}
\label{sec:09_decays}

\textbf{What it does.} Pythia8 decays every unstable hadron whose mean decay length $c\tau$ is below $10\mm$ and
leaves the others to the detector simulation; the $B$ and $D$ mesons of our event ($c\tau$ of $0.5$ and
$0.12\mm$) are decayed by the generator, the $K^0_S$ ($27\mm$), the $\Lambda$ ($79\mm$), the charged pions and
kaons (metres) are not. A decayed hadron stays in the record with status 2 and a decay vertex. The
\code{Generator} then asks where that vertex is. Inside the $2.9\cm$ cylinder around the beam, the decay
happened in the vacuum of the beam pipe and is kept as Pythia made it: Geant4 gets the daughters as primaries,
starting at the decay point, which is stored as a \code{SimVertex} without a parent track. Outside the cylinder the
parent itself is given to Geant4, with Pythia's decay attached as a ``predefined decay channel'', so that it decays
where the generator decided but inside the material, after having crossed it.

\textbf{The parameters.} \code{ParticleDecays:limitTau0 = on}, \code{ParticleDecays:tau0Max = 10} (in mm,
\path{Configuration/Generator/python/Pythia8CommonSettings_cfi.py}); \code{RDecLenCut = 2.9} (cm) in the
\code{Generator} parameter set; its partner \code{LDecLenCut = 30.0}, a length in $z$, is read and never used
(section~5).

\textbf{The numbers of the event.} The file holds 29 \code{SimVertex}s for the signal although nothing decayed in
Geant4: the smeared primary vertex and 28 generator decay vertices, all inside the cylinder, so every daughter of
the event starts as a primary. Table~\ref{tab:09_decays} follows the three weak decays of the two $b$ systems. The
distance a particle flies before it decays is a random number, exponential with mean $\beta\gamma c\tau$, and the
three in the table illustrate it: the $B^+$ flew half of its mean, the $\bar D^0$ more than twice its own, which
happens one time in ten. The $\bar D^0$ of the $\bar b$ jet decayed $25.8\cm$ from the collision, still inside the
beam pipe at $r = 1.2\cm$, because a particle at $\eta=-3.76$ moves $21$ times faster along the beam than away from
it. The primaries of these decays start from those points (figure~\ref{fig:09_rz}).

\begin{table}[!htb]
\centering\small\setlength{\tabcolsep}{4pt}
\caption{The three weak decays of the $b$ and $\bar b$ systems, kept from the generator (FWLite extras of the
companion; $m$ and $c\tau$ from the Particle Data Group: $B^\pm$ $5.279\GeV$, $491\um$; $D^0$ $1.865\GeV$,
$123\um$). The flight is measured from the point where the hadron was made.}
\label{tab:09_decays}
\begin{tabular}{@{}l rrr r r rr l@{}}
\toprule
hadron & $E$ & \pt{} & $\eta$ & $\beta\gamma$ & mean $\beta\gamma c\tau$ & flight & decay $(r,z)$ & daughters \\
 & [GeV] & [GeV] & & & [cm] & [cm] & [cm] & \\
\midrule
$B^-$ ($b$ jet) & 235.9 & 94.3 & $-1.57$ & 44.7 & 2.2 & 0.78 & $(0.32, 2.39)$ & $\bar\nu_e$, $e^-$, $D_0^*(2300)^0$ \\
$B^+$ ($\bar b$ jet) & 2388.2 & 112.6 & $-3.75$ & 452 & 22.2 & 11.7 & $(0.55, -8.60)$ & $\bar D^{*0}$, $\pi^+$, $\pi^0$, $K^0$, $\bar K^{*0}$ \\
$\bar D^0$ (from the $B^+$) & 911.9 & 42.4 & $-3.76$ & 489 & 6.0 & 14.1 & $(1.21, -22.66)$ & $\bar K^0$, $\pi^-$, $\pi^+$, $\pi^0$, $\pi^0$ \\
\bottomrule
\end{tabular}
\end{table}

\subsection{Straight lines to the calorimeters}
\label{sec:09_boundary}

\textbf{What it does.} With no process registered, the only thing Geant4 does to a particle is the
\emph{transportation}: it moves it in a straight line to the next boundary between two volumes, then to the next,
until the particle leaves the world or reaches the time limit of the \code{SteppingAction}
(\code{MaxTrackTime = 500}\,ns). The step that takes a track from an ``inside'' volume into a calorimeter
volume is seen by \code{CaloTrkProcessing}, which writes the point and the momentum into the \code{SimTrack} (\code{trackerSurfacePosition()}, \code{trackerSurfaceMomentum()}, \code{crossedBoundary()}).
This is the moment at which the truth of a particle passes from the tracker to the calorimeters: in the next
chapters a shower in the ECAL is credited to the track that crossed here, not to the electrons it is made of.

\textbf{The parameters.} The surfaces are those of figure~\ref{fig:09_rz}. Two numbers follow from them by
geometry. Seen from our vertex at $z=+3.11\cm$, the corner of the tracker volume, $(r,z)=(123.3,-293.2)\cm$, is at
$\eta=-\sinh^{-1}(296.3/123.3) = -1.61$: on the negative side a particle leaves the tracker through the barrel
cylinder if $\abseta<1.61$, through the endcap disk if it is beyond. And the hole of the HF closes the acceptance:
a line from the centre reaches the inner radius of the HF, $12.5\cm$, at its back face, $|z|=1280\cm$, if
$\sinh\abseta = 1280/12.5 = 102.4$, that is $\abseta = 5.32$. A particle beyond that passes through the hole of
every calorimeter.

\textbf{The numbers of the event.} 557 of the 571 primaries cross into a calorimeter volume, with exactly their
production momentum: no field, no material effect. Table~\ref{tab:09_boundary} says through which surface. The $b$
system, at $\eta=-1.55$, is $0.06$ from the corner: 7 of its 8 primaries leave the tracker through the barrel
cylinder between $z=-240$ and $-292\cm$, and one photon of $1.5\GeV$ at $\eta=-1.68$ through the endcap disk at
$r=114\cm$. The $b$ jet that the calorimeters and the trigger of chapters~\ref{ch:11_DIGI} and~\ref{ch:12_L1T} find on
the barrel--endcap boundary sits there already in the geometry of the tracker. The $\bar b$ system crosses, all 16, through the ring at the end of the nose,
$z=-355.1\cm$, at $r=15.4$--$20.3\cm$, and flies another $7.6$\,m to the front face of the HF. The 14 primaries
that never cross are 13 particles at $5.34<\abseta<5.50$, accepted by the $\eta$ cut but beyond the $5.32$ of the
HF hole, and the forward neutron. The usable acceptance of the simulation therefore ends at $\abseta\approx5.3$,
not at the $5.5$ of the cut; the 16 signal particles with $5.2<\abseta<5.5$ ($836\GeV$) are all accepted and only
three of them reach a calorimeter volume.

\begin{table}[!htb]
\centering\small
\caption{Where the primaries cross into a calorimeter volume (\code{extras\_09\_G4.txt}, from
\code{SimTrack::trackerSurfacePosition}), for the signal and for library event 72; the $\abseta$ range is that of
the signal primaries through the surface.}
\label{tab:09_boundary}
\begin{tabular}{@{}lll rr@{}}
\toprule
surface & where & $\abseta$ (signal) & signal & event 72 \\
\midrule
barrel cylinder & $r=123.3\cm$, $|z|<293.2\cm$ & 0.01--1.61 & 172 & 54 \\
endcap disk & $|z|=293.2\cm$, $r=23$--$122\cm$ & 1.60--3.23 & 228 & 66 \\
end of the nose & $|z|=355.1\cm$, $r=7.8$--$21.3\cm$ & 3.50--4.51 & 98 & 18 \\
front face of the HF & $|z|=1114.9\cm$, $r=13$--$85\cm$ & 3.26--5.16 & 36 & 9 \\
inner walls along the beam pipe & $r=7.8$--$12.5\cm$, $|z|=356$--$1277\cm$ & 4.51--5.32 & 23 & 8 \\
\midrule
no crossing & & 5.34--5.50, 9.72 & 14 & 9 \\
all primaries & & & 571 & 164 \\
\bottomrule
\end{tabular}
\end{table}

\begin{figure}[!htb]
\centering
\begin{subfigure}[t]{0.495\textwidth}
\includegraphics[width=\textwidth]{09_G4/out/simAnalyzer_signal_09_G4_page-07.png}
\caption{the 557 primaries at the boundary, marker at the crossing point: page~6 again.}
\end{subfigure}\hfill
\begin{subfigure}[t]{0.495\textwidth}
\includegraphics[width=\textwidth]{09_G4/out/simAnalyzer_signal_09_G4_page-08.png}
\caption{the secondaries created by Geant4: none, only the primaries as grey dots.}
\end{subfigure}
\caption{The signal through the transport-only Geant4, pages 7 and 8 of \code{simAnalyzer\_signal\_09\_G4.pdf}
(\code{\_page-07.png}, \code{\_page-08.png}). Without field and physics every primary arrives where it was
aimed and creates nothing: the two pages that will change most in chapter~\ref{ch:10_SIM}.}
\label{fig:09_boundary}
\end{figure}

\subsection{Hits without physics}
\label{sec:09_hits}

\textbf{What it does.} A sensitive detector is the code that a sensitive volume calls at every step taken inside
it. The tracker one (\code{TkAccumulatingSensitiveDetector}) makes one \code{PSimHit} per crossing of a silicon
sensor: entry and exit point in the frame of the sensor, energy loss, time of flight, the track that made it.
Normally a crossing with no energy loss is not recorded; here \code{TrackerSD.ZeroEnergyLoss = True} records it
anyway, so that the transport-only run shows which sensors each straight line crosses. The hits are split by their
time of flight into a \code{LowTof} and a \code{HighTof} collection, the in-time and the late ones, at three times
the shaping time of the strip electronics. The calorimeter one (\code{CaloSD} and its daughters \code{ECalSD},
\code{HCalSD}) adds the energy deposited in a cell by a track into one \code{PCaloHit} per cell, track and time
slice; with no physics nothing is deposited, except in the HF (box below).

\begin{plainwords}[The HF shower library]
The HF is steel with quartz fibres in it, and it sees the Cherenkov light that the charged particles of a shower
make in the fibres: far too many photons to follow one by one in every event. CMS simulated such showers once,
electromagnetic and hadronic ones at a set of energies, and stored the light captured in the fibres (the file
\code{HFShowerLibrary\_run3\_v7.root}). When a photon, electron or hadron (not a muon) takes its first step in
the HF, \code{HCalSD} (\code{getFromHFLibrary}) takes one stored shower from each of the two library energies
around the particle's energy, keeps their photons in proportion, turns them to the particle's direction and
position, applies the attenuation along the fibre, books each surviving photon as a photoelectron, a hit of
``$1\GeV$'' in its cell, and kills the particle. The library is called by the sensitive detector, not by a
physics process, which is why the HF answers even here.
\end{plainwords}

\textbf{The parameters.} \code{TrackerSD}: \code{ZeroEnergyLoss = False} by default (\code{True} by the customise
line of this chapter), \code{ElectronicSigmaInNanoSeconds = 12.06}, hence the \code{LowTof}/\code{HighTof} split at
$3\times12.06 = 36.2\ns$, \code{EnergyThresholdForPersistencyInGeV = 0.2}. \code{CaloSD}: \code{EminHits = 15} and
$10\MeV$ for EB and EE, \code{TmaxHits = 500}\,ns. \code{HFShower}: \code{UseShowerLibrary = True}, the library
file set by the modifier \code{run3\_HFSL} of \path{Geometry/HcalSimData/python/HFParameters_cff.py}.

\textbf{The numbers of the event.} The tracker records 6734 \code{PSimHit}s (PXB 908, PXF 829, TIB 960, TID
679, TOB 860, TEC 2498), all in the \code{LowTof} collections (time of flight $0.2$ to $14.0\ns$), from 375
primaries, 193 charged and 182 neutral: a photon ``hits'' a sensor here because it crosses it. The count is pure
geometry. A straight track at $\abseta<1$ crosses the 4 pixel layers, the 4 layers of the inner barrel (TIB) and the 6 of the outer
barrel (TOB); the first two layers of TIB and of TOB carry two sensors back to back, one with its strips at a
small stereo angle, so the track crosses $4+6+8=18$ sensors. The 38 charged primaries with $\abseta<1$ have $18.7$
hits on average (17 to 21: neighbouring modules overlap a little so as to leave no gap), the 47 photons $18.6$.
Between $\abseta=1$ and 2, where a track crosses barrel layers and then endcap disks, the average is $21.3$ and
the most, 32, belongs to a photon at $\eta=-1.96$. The $b$ system leaves 179 hits; the $\bar b$ system, at
$\eta=-3.75$, none.

Of the calorimeters only the HF responds: 746 \code{PCaloHit}s in 459 cells, 1511 photoelectrons, from 144
primaries that carried $4266\GeV$ into it, $0.35$ photoelectrons per GeV. A shower is not an average: the
$\pi^+$ of $643\GeV$ of the $\bar b$ system gives 321 photoelectrons ($0.50$ per GeV), a $K^0_L$ of $241\GeV$ next
to it 84 ($0.35$), a $\pi^+$ of $165\GeV$ 46 ($0.28$). The $\bar b$ system puts 760 photoelectrons into 39 cells.
Figure~\ref{fig:09_hits} shows both pages.

\begin{figure}[htb]
\centering
\begin{subfigure}[t]{0.495\textwidth}
\includegraphics[width=\textwidth]{09_G4/out/simAnalyzer_signal_09_G4_page-09.png}
\caption{the 6734 tracker \code{PSimHit}s, zero energy loss, by subdetector; red: the $b$ system.}
\end{subfigure}\hfill
\begin{subfigure}[t]{0.495\textwidth}
\includegraphics[width=\textwidth]{09_G4/out/simAnalyzer_signal_09_G4_page-10.png}
\caption{the \code{PCaloHit}s: only the 459 HF cells of the shower library; blue: the $\bar b$ system.}
\end{subfigure}
\caption{The signal through the transport-only Geant4: the tracker hits (left) and the calorimeter hits (right),
pages 9 and 10 of \code{simAnalyzer\_signal\_09\_G4.pdf} (\code{\_page-09.png}, \code{\_page-10.png}). Pages 1--5
of the same PDF are the five generator stages, pages 6--8 the primaries, the boundary crossing and the (empty) list
of secondaries.}
\label{fig:09_hits}
\end{figure}


\subsection{The pileup collisions}
\label{sec:09_pileup}

\textbf{What it does.} The library of chapter~\ref{ch:08_PU} goes through the same transport-only Geant4, all 250
collisions, each from its own smeared vertex. Nothing is mixed yet: each collision is an event of its own, and
the ledger sums the 68 that the mixing of chapter~\ref{ch:10_SIM} will put into our bunch crossing.

\textbf{The parameters.} Those of the signal, with the Geant4 seed of the library, \code{PUSEED}$+2 = 20243$
(only the HF shower library draws random numbers in this run).

\textbf{The numbers of the event.} Library event~72, the first in-time collision (figure~\ref{fig:09_pileup}), has
192 stable particles of which 164 become primaries; 27 are beyond $\abseta=5.5$ with $9041$ of its $13\,600\GeV$
(three of them beyond $8.2$, among them a $\Lambda$ of $3414\GeV$ at $\eta=9.85$), one photon is below $40\MeV$,
and four forward neutrals with $3749\GeV$ are kept by the very-forward rule. 155 of the 164 reach the
calorimeter boundary (table~\ref{tab:09_boundary}: 54 through the barrel cylinder, 66 through the endcap disk);
the collision leaves 1965 tracker hits and 153 HF hits with 192 photoelectrons in 135 cells. Over the 68
collisions the number of primaries runs from 0 (library event 124, a diffractive collision with 6 stable
particles, all beyond $\abseta=5.5$) to 470 (event 94, 547 stable particles), 153 on average.

\begin{figure}[!htb]
\centering
\begin{subfigure}[t]{0.495\textwidth}
\includegraphics[width=\textwidth]{09_G4/out/simAnalyzer_pileup_09_G4_page-06.png}
\caption{the 164 primaries and the 28 stable particles without a \code{SimTrack} (grey crosses).}
\end{subfigure}\hfill
\begin{subfigure}[t]{0.495\textwidth}
\includegraphics[width=\textwidth]{09_G4/out/simAnalyzer_pileup_09_G4_page-09.png}
\caption{its 1965 tracker \code{PSimHit}s by subdetector, zero energy loss.}
\end{subfigure}
\caption{One pileup collision through the transport-only Geant4: library event~72, pages 6 and 9 of
\code{simAnalyzer\_pileup\_09\_G4.pdf} (\code{\_page-06.png}, \code{\_page-09.png}). No jet, no $b$ system: an even
spray of soft particles.}
\label{fig:09_pileup}
\end{figure}

\section{Ledger and display}

Table~\ref{tab:09_ledger} is the ledger of the hand-over, by region and by kind, for the signal, for the 68
collisions and for their sum. For the signal the stage removes 63 of 634 particles, but $6677\GeV$ of
$13\,600$: 90\,\% of the particles, 89\,\% of the \pt{} and 51\,\% of the energy go into Geant4. Everything in
$2.5<\abseta<5.5$ passes untouched; the central region loses the six soft photons and the neutrino; beyond
$5.5$ only the forward neutron and the border $\pi^+$ survive. For the pileup the picture is more extreme:
80\,\% of the particles and 84\,\% of the \pt{} enter, but only 24\,\% of the energy, because $840$ of the
$925\TeV$ of a bunch crossing's minimum-bias collisions sit beyond $\abseta=5.5$ in the beam remnants and
$706\TeV$ of them are never simulated. What remains per collision is on average 153 primaries, $3.2\TeV$
and $66\GeV$ of \pt{}, of which 70 particles and $32\GeV$ of \pt{} in the tracker region. The sum of signal and
68 collisions is what chapter~\ref{ch:10_SIM} will mix: $10\,959$ primaries with $225\TeV$ and $4957\GeV$ of
\pt{}, the signal being 5\,\% of the particles and 10\,\% of the \pt{} of its own bunch crossing.

\begin{table}[htbp]
\centering\scriptsize\setlength{\tabcolsep}{2.5pt}
\caption{The ledger of the hand-over (\code{ledger\_09\_G4.txt}; the regions and kinds from the FWLite extras of
the companion): the stable particles before and the Geant4 primaries after, $n$ / $E$ [GeV] / \pt{} [GeV],
for the signal and for the 68 in-time pileup collisions; the last block gives the fractions handed over, the dropped
particles by the reasons of table~\ref{tab:09_cuts}, and the sum.}
\label{tab:09_ledger}
\begin{tabular}{lrrrrrrcrrrrrr}
\toprule
 & \multicolumn{6}{c}{signal} & & \multicolumn{6}{c}{68 pileup collisions} \\
\cmidrule(lr){2-7}\cmidrule(lr){9-14}
 & \multicolumn{3}{c}{before: stable} & \multicolumn{3}{c}{after: primaries} & & \multicolumn{3}{c}{before: stable} & \multicolumn{3}{c}{after: primaries} \\
 & $n$ & $E$ & \pt{} & $n$ & $E$ & \pt{} & & $n$ & $E$ & \pt{} & $n$ & $E$ & \pt{} \\
\midrule
$\abseta<2.5$     & 316 & 733.8 & 278.7 & 309 & 638.1 & 241.1 & & 4864 & 5519 & 2179.4 & 4765 & 5506 & 2173.6 \\
$2.5<\abseta<3.0$ & 55 & 193.1 & 24.6 & 55 & 193.1 & 24.6 & & 1037 & 3597 & 449.2 & 1035 & 3597 & 449.2 \\
$3.0<\abseta<5.2$ & 189 & 5019.6 & 197.4 & 189 & 5019.6 & 197.4 & & 3996 & 58\,413 & 1674.3 & 3992 & 58\,367 & 1672.7 \\
$5.2<\abseta<5.5$ & 16 & 836.4 & 7.8 & 16 & 836.4 & 7.8 & & 436 & 17\,054 & 160.3 & 428 & 16\,888 & 158.9 \\
$\abseta>5.5$     & 58 & 6817.1 & 20.3 & 2 & 235.9 & 0.7 & & 2660 & 840\,217 & 851.4 & 168 & 134\,196 & 31.1 \\
\midrule
charged hadrons & 303 & 7167.7 & 259.1 & 276 & 3298.6 & 247.8 & & 5911 & 572\,735 & 3087.8 & 4677 & 55\,664 & 2596.3 \\
neutral hadrons & 56 & 3449.9 & 84.7 & 48 & 1869.8 & 81.1 & & 1088 & 179\,328 & 749.7 & 837 & 80\,592 & 627.3 \\
photons         & 273 & 2875.9 & 142.7 & 246 & 1743.7 & 137.9 & & 5900 & 171\,676 & 1458.8 & 4808 & 82\,175 & 1251.9 \\
electrons, muons & 1 & 10.9 & 5.0 & 1 & 10.9 & 5.0 & & 87 & 1004 & 12.6 & 66 & 121.4 & 10.1 \\
neutrinos       & 1 & 95.5 & 37.5 & 0 & 0 & 0 & & 7 & 57.5 & 5.6 & 0 & 0 & 0 \\
\midrule
all & 634 & 13\,600.0 & 528.9 & 571 & 6923.0 & 471.7 & & 12\,993 & 924\,800 & 5314.6 & 10\,388 & 218\,553 & 4485.6 \\
\bottomrule
\end{tabular}

\medskip
\begin{tabular}{ll}
\toprule
signal & handed over: 90.1\,\% of the particles, 50.9\,\% of $E$, 89.2\,\% of \pt{}; \\
 & dropped ($E$ in GeV): 56 beyond $\abseta=5.5$ (6581.3), 6 below $40\MeV$ (0.2), 1 neutrino (95.5) \\
68 collisions & handed over: 80.0\,\% of the particles, 23.6\,\% of $E$, 84.4\,\% of \pt{}; \\
 & dropped: 2492 beyond $\abseta=5.5$ ($706\,021$), 98 below $40\MeV$ (2.7), 7 neutrinos (57.5), 8 other (166.4) \\
signal + pileup & before: $13\,627$ stable particles, $938\,400\GeV$, $5843.5\GeV$ of \pt{}; \\
 & after: $10\,959$ primaries, $225\,476\GeV$, $4957.3\GeV$ of \pt{} (80.4\,\%, 24.0\,\%, 84.8\,\%); \\
 & dropped: 2548 beyond $\abseta=5.5$ ($712\,602$), 104 below $40\MeV$ (2.9), 8 neutrinos (153.0), 8 other (166.4) \\
\bottomrule
\end{tabular}
\end{table}

Figure~\ref{fig:09_prim} is the $\eta$--$\phi$ display of the event at this stage: the 571 primaries in the
colour of their kind, marker size growing with \pt{}, the $b$ system in red at $(\eta,\phi) = (-1.55, -2.09)$
and the $\bar b$ system in blue at $(-3.75, 1.01)$, and the 63 particles without a \code{SimTrack} as grey
crosses. The crosses are where the ledger says: a column at each border of the plot (38 of the 56 forward
particles sit beyond the $|\eta|=6$ edge of the display and are drawn on it), six faint ones scattered over the
centre (the photons below $40\MeV$) and one at the $b$ system (the neutrino). Compared with page~5 of the
same PDF, the stage-5 display of chapter~\ref{ch:06_Decay}, nothing has moved: the hand-over is a selection,
not a transformation, and page~7 (figure~\ref{fig:09_boundary}) is the same picture again.

\begin{figure}[htbp]
\centering
\includegraphics[width=0.7\textwidth]{09_G4/out/simAnalyzer_signal_09_G4_page-06.png}
\caption{The signal handed to Geant4 (\code{simAnalyzer\_signal\_09\_G4\_page-06.png}): the 571 primaries
(\code{SimTrack}s with a generator link) by kind, the $b$ and $\bar b$ systems in red and blue with their
vector sum (star), and the 63 stable particles without a \code{SimTrack} as grey crosses (56 beyond
$\abseta=5.5$, 6 below $40\MeV$, 1 neutrino). The right panel marks the SIM stage of this sample: transport
only.}
\label{fig:09_prim}
\end{figure}

\begin{plainwords}[What the pileup brings into Geant4]
A minimum-bias collision carries $13.6\TeV$, the same as the signal, but almost all of it goes down the beam
pipe: in the 68 collisions $840\TeV$ of $925$ are beyond $\abseta=5.5$, where there is no detector, and the
hand-over drops $706\TeV$ of them before Geant4 does a single step. What enters is small per collision, 153
particles and $66\GeV$ of \pt{}, but there are 68 of them: $4486\GeV$ of \pt{} against the $472$ of the
signal. In the tracker region alone the pileup primaries carry $2174\GeV$ of \pt{} to the signal's $241$. In the
tracker region the 68
collisions lay down $2174/(5\times2\pi) = 69\GeV$ of \pt{} per unit of $\eta$--$\phi$ area; at the HL-LHC, with
200 collisions, the same hand-over would deliver about $200\GeV$.
\end{plainwords}

\FloatBarrier
\section{Open problems for Phase-2}

The hand-over is configured once for all eras in \path{SimG4Core/Application/python/g4SimHits_cfi.py}. A look
through that file and the Phase-2 modifiers it imports gives the following.

\textbf{The $\eta$ cut is not era-dependent.} \code{Generator.MinEtaCut = -5.5}, \code{MaxEtaCut = 5.5} and
\code{RDecLenCut = 2.9} carry no modifier at all; the only places where the cut changes are
\path{Configuration/StandardSequences/python/SimWithCastor_cff.py} ($-7.0$ to $5.5$, for the CASTOR years) and
\path{SimG4Core/Application/python/NeutronBGforMuons_cff.py} ($\pm7.0$, a neutron-background study). The Phase-2
detector list of \code{phase2\_common} in the same file (\code{OnlySDs}) adds the HGCal, the MTD
(\code{MtdSensitiveDetector}), the HFNose and keeps the ZDC, the BCM1F and the PLT; the Phase-2 Level-1 regions of
\path{L1Trigger/Phase2L1ParticleFlow/python/l1ctLayer1_cff.py} put the HGCal at $1.5<\abseta<3.0$ and the HF at
$3.0<\abseta<5.5$. In this event the particles with $5.2<\abseta<5.5$ (16, $836\GeV$) are accepted and 13 of
them never enter a calorimeter volume (the hole of the HF, section~\ref{sec:09_boundary}); whether the
Phase-2 geometry closes that gap could not be established from the python. The very-forward rule that keeps neutrons and photons aimed at the ZDC
(\code{EtaCutForHector = 8.2}, 168 primaries with $134\TeV$ in our 68 collisions) runs at the HL-LHC as it does
now, with \code{SteppingAction.CMStoZDCtransport = False}, so these particles are simulated up to the beam-line
dead regions and then killed: CPU spent on tracks that produce nothing.

\textbf{The $40\MeV$ cut and the decay cylinder.} \code{MinPCut = 0.04} has no modifier either. It is a cut
on the total momentum, below the response of any detector of Run~3; for the HGCal (\code{HGCSD.EminHit = 0.0})
and the MTD (\code{MtdSD.EminHit = 0.0}) the sensitive detectors keep every deposit, so the $40\MeV$ remains the
only lower limit of what is simulated, and it has not been revisited. \code{LDecLenCut = 30.0} is read in
\code{Generator.cc} into a local variable, printed and never used: the decay cylinder is defined by its radius
alone, $2.9\cm$, the Run-3 beam pipe. The $\bar D^0$ of this event decayed at $r = 1.2\cm$, $z = -22.7\cm$,
inside it, so Pythia's decay was kept; a $2.9\cm$ cylinder of unlimited length also keeps generator decays that
happen metres down the beam pipe. The radius of the Phase-2 beam pipe and of the first inner-tracker layer is
not in a python file of the release and I could not establish whether $2.9\cm$ still corresponds to it.

\textbf{The StackingAction thresholds.} The Russian roulette of the \code{StackingAction} kills low-energy
photons (\code{RusRoGammaEnergyLimit = 5} MeV) and neutrons (\code{RusRoNeutronEnergyLimit = 10} MeV) with a
survival probability per region: $0.3$ for photons and $0.1$ for neutrons in the ECAL, HCAL, muon iron,
preshower and world. The two Phase-2 entries, \code{RusRoHGcalGamma = 1.3} and \code{RusRoHGcalNeutron = 1.1},
are not tuned probabilities: \path{SimG4Core/Application/src/StackingAction.cc} applies the roulette only when
the value is below 1 (line 341, \code{prob < 1.0}) and only registers the \code{HGCalRegion} in that case
(line 399), so these values switch the roulette off in the HGCal. Every photon above the production cut and
every neutron is tracked there, which is the right choice for a calorimeter with $0.5$ and $1$\,cm$^2$ cells
and timing, and the CPU cost of it at PU 200 is paid in full. The heavy-particle thresholds
(\code{NeutronThreshold}, \code{ProtonThreshold}, \code{IonThreshold} $=30\MeV$) act only with
\code{KillHeavy = True}, which is off.

\textbf{Physics list and Geant4 version.} The signal log carries the banner ``Geant4 version Name:
geant4-11-02-patch-02 [MT] (21-June-2024)'', i.e.\ the \code{geant4} tool 11.2.2 of this release. The list is
\code{SimG4Core/Physics/FTFP\_BERT\_EMM} for every era (no modifier in \code{g4SimHits\_cfi.py}; only the
CPU-benchmark and neutron-background customisations switch to \code{FTFP\_BERT}, \code{FTFP\_BERT\_XS\_EML},
\code{FTFP\_BERT\_HP\_EML}): \path{SimG4Core/PhysicsLists/plugins/FTFPCMS_BERT_EMM.cc} registers the CMS
electromagnetic constructor \code{CMSEmStandardPhysics}, \code{G4DecayPhysics}, elastic, FTFP/Bertini
(\code{EminFTFP = 3}, \code{EmaxBERT = 6}, \code{EmaxBERTpi = 12} GeV) and stopping physics. The EM options
\code{G4TransportWithMSC = 2}, \code{G4GammaGeneralProcess = True}, \code{G4HepEmActive = False} are the same
for Phase-2. Whether the HGCal needs a different EM constructor (the \code{\_EMMT} and \code{\_EMM\_TRK}
variants exist next to it) could not be established from the configuration.

\textbf{The time windows.} \code{common\_maximum\_time} sets \code{MaxTrackTime = 500} ns in the central region,
\code{MaxTrackTimeForward = 2000} ns beyond \code{MaxZCentralCMS = 50} m, and the only modifier that touches it
is \code{h2tb} (a test beam, 1000\,ns); \code{CaloSD.TmaxHits} is $500\ns$ for every calorimeter collection
and $2000$ for the ZDC. Phase-2 changes the stepping, not the window: \path{SimG4Core/Application/src/RunManagerMT.cc}
(\code{runForPhase2}) switches to \code{Phase2SteppingAction} as soon as a region named \code{HGCalRegion},
\code{FastTimerRegionBTL} or \code{FastTimerRegionETL} exists in the geometry, and that class records the
tracker-to-BTL transition (\code{BTLName = BarrelTimingLayer}) with the same time cuts. The timing detectors are
given finer slices (\code{MtdSD.TimeSliceUnit = 0.01} ns, \code{HGCSD.TimeSliceUnit = 0.001} ns, and
\code{phase2\_timing} sets \code{ECalSD.TimeSliceUnit = 0.001} with \code{StoreLayerTimeSim = True}), while
\code{MtdSD.EnergyThresholdForPersistencyInGeV = 1e9} and \code{ForHistoryInGeV = 1e9} are marked in the file
``keep temporarily old behaviour'': no MTD hit makes its track persistent.

\textbf{What \code{phase2} modifies in \code{g4SimHits\_cfi.py}, and what it does not.} The modifiers found by
grep are \code{phase2\_hgcal} (\code{common\_MCtruth.SaveCaloBoundaryInformation = True}), \code{phase2\_timing}
(the ECAL time slices above), \code{phase2\_common} (the \code{OnlySDs}, \code{TrackHits} and \code{CaloHits}
lists with \code{FastTimerHitsBarrel/Endcap}, \code{HGCHitsEE/HEfront/HEback}, \code{HFNoseHits};
\code{LHCTransport = False}, i.e.\ no PPS proton transport in Phase-2 where \code{run3\_common} sets it
\code{True}; in \code{MuonSD}, \code{HaveDemoChambers = False}), \code{phase2\_hgcalOnly} and
\code{phase2\_hgcalV18} (in \code{HGCSD}, \code{HitCollection = 2}). Nothing is era-modified in
\code{Generator}, in \code{StackingAction}, in \code{Physics}, in \code{common\_maximum\_time} or in
\code{TrackerSD}, whose \code{EnergyThresholdForPersistencyInGeV = 0.2} and
\code{ElectronicSigmaInNanoSeconds = 12.06} (hence the $36.2\ns$ Low/High time-of-flight split, designed for
the strip electronics) are inherited by the Phase-2 pixel modules: the hand-over of the HL-LHC is,
parameter for parameter, the hand-over of this chapter.
